% OmniLaTeX User Manual
% Compile with: latexmk -lualatex omnilatex.tex
%
% This is the CTAN documentation for the omnilatex package.

\documentclass[
  doctype=manual,
  language=english,
]{omnilatex}

% ── Manual-specific setup ────────────────────────────
\usepackage{bookmark}
\usepackage{multicol}
\usepackage{longtable}

\hypersetup{
  pdftitle={OmniLaTeX User Manual},
  pdfauthor={Wyatt Au},
}

% ── Helper commands ──────────────────────────────────
\newcommand{\option}[1]{\texttt{#1}}
\newcommand{\module}[1]{\texttt{#1}}
\newcommand{\file}[1]{\texttt{#1}}

\title{OmniLaTeX}
\subtitle{A Modular Document Class for Academic and Professional Documents}
\author{Wyatt Au}
\date{Version 1.0.0 — 2026-04-03}

\begin{document}

\maketitle
\tableofcontents

%% ─────────────────────────────────────────────────────
\chapter{Introduction}
%% ─────────────────────────────────────────────────────

OmniLaTeX is a modular LaTeX document class built on LuaLaTeX and KOMA-Script.
It provides 46 doctype aliases, 21 modular packages, and support for
academic and professional documents including theses, articles, CVs,
technical reports, patents, and more.

\section{Requirements}

\begin{itemize}
	\item LuaLaTeX (TeX Live 2024 or newer)
	\item KOMA-Script (bundled with TeX Live)
	\item polyglossia (for multilingual support)
\end{itemize}

\section{Installation}

\subsection{From CTAN (recommended)}

Once published on CTAN, install via:

\begin{minted}{bash}
tlmgr install omnilatex
\end{minted}

\subsection{From Source}

Clone the repository and ensure the \file{lib/} and \file{config/}
directories are in your \texttt{TEXINPUTS} path.

\section{Quick Start}

\begin{minted}{latex}
\documentclass[doctype=article]{omnilatex}
\begin{document}
\title{My Document}
\author{Author Name}
\maketitle
\tableofcontents
\section{Introduction}
Hello, world!
\end{document}
\end{minted}

Compile with:

\begin{minted}{bash}
latexmk -lualatex main.tex
\end{minted}

%% ─────────────────────────────────────────────────────
\chapter{Class Options}
%% ─────────────────────────────────────────────────────

OmniLaTeX accepts the following options via \texttt{\textbackslash documentclass}:

\begin{longtable}{lp{8cm}}
	\toprule
	\textbf{Option}            & \textbf{Description}                                                                   \\
	\midrule
	\option{language}          & Document language (default: \option{english}). Passed to polyglossia.                  \\
	\option{doctype}           & Document type alias (default: \option{thesis}). See Chapter~\ref{ch:doctypes}.         \\
	\option{titlestyle}        & Title page style (default: \option{book}).                                             \\
	\option{institution}       & Institution configuration (default: \option{none}). See Chapter~\ref{ch:institutions}. \\
	\option{censoring}         & Enable redaction support (default: \option{false}).                                    \\
	\option{loadGlossaries}    & Enable glossaries module (default: \option{false}).                                    \\
	\option{todonotes}         & Enable TODO notes (default: \option{false}).                                           \\
	\option{enablefonts}       & Enable custom font loading (default: \option{false}).                                  \\
	\option{enablegraphics}    & Enable graphics module (default: \option{false}).                                      \\
	\option{enablemath}        & Enable math module (default: \option{true}).                                           \\
	\option{enabletikz}        & Enable TikZ modules (default: \option{true}).                                          \\
	\option{enableengineering} & Enable engineering TikZ shapes (default: \option{true}).                               \\
	\option{enablecode}        & Enable code listings (default: \option{true}).                                         \\
	\option{enabletables}      & Enable table module (default: \option{true}).                                          \\
	\bottomrule
\end{longtable}

%% ─────────────────────────────────────────────────────
\chapter{Document Types}
\label{ch:doctypes}
%% ─────────────────────────────────────────────────────

OmniLaTeX provides 46 doctype aliases that resolve to 13 document profiles
across 3 KOMA-Script base classes.

\section{Available Doctypes}

\begin{longtable}{lll}
	\toprule
	\textbf{Alias}  & \textbf{Profile} & \textbf{Base Class} \\
	\midrule
	thesis          & Thesis           & scrbook             \\
	dissertation    & Dissertation     & scrbook             \\
	book            & Book             & scrbook             \\
	manual          & Manual           & scrbook             \\
	article         & Article          & scrartcl            \\
	journal         & Journal          & scrartcl            \\
	cv              & CV               & scrartcl            \\
	cover-letter    & Cover Letter     & scrartcl            \\
	technicalreport & Technical Report & scrreprt            \\
	standard        & Standard         & scrreprt            \\
	patent          & Patent           & scrreprt            \\
	dictionary      & Dictionary       & scrbook             \\
	inlinepaper     & Inline Paper     & scrartcl            \\
	\bottomrule
\end{longtable}

Additional aliases exist for common variations (e.g., \option{report},
\option{paper}, \option{resume}, \option{letter}).
See \file{omnilatex.cls} for the complete alias table.

\section{Doctype Profiles}

Each doctype profile (in \file{config/document-types/}) defines:

\begin{itemize}
	\item Page layout (\texttt{DIV}, margins, headers/footers)
	\item Font size
	\item Line spacing
	\item Title page style
	\item Section numbering depth
	\item Table of contents depth
	\item Color mode (color, grayscale, bw)
\end{itemize}

%% ─────────────────────────────────────────────────────
\chapter{Institution Configuration}
\label{ch:institutions}
%% ─────────────────────────────────────────────────────

OmniLaTeX supports pluggable institution branding via
\file{config/institutions/}.

\section{Using an Institution}

\begin{minted}{latex}
\documentclass[doctype=thesis,institution=tuhh]{omnilatex}
\end{minted}

This loads \file{config/institutions/tuhh/tuhh.sty}, which provides:

\begin{itemize}
	\item Logo definitions (language-aware via \texttt{\textbackslash DeclareTranslation})
	\item Institution links
	\item Color palette
	\item Custom commands (degree names, department names, etc.)
\end{itemize}

\section{Creating a Custom Institution}

See \texttt{CONTRIBUTING.md} for a step-by-step guide.

%% ─────────────────────────────────────────────────────
\chapter{Modules}
\label{ch:modules}
%% ─────────────────────────────────────────────────────

OmniLaTeX is organized into 21 modules across 9 subdirectories:

\section{Core (\module{lib/core/})}

\begin{description}
	\item[\module{omnilatex-base}] Build mode detection (dev/prod/ultra), date/time utilities.
\end{description}

\section{Layout (\module{lib/layout/})}

\begin{description}
	\item[\module{omnilatex-page}] Headers, footers, page makeup via scrlayer-scrpage.
	\item[\module{omnilatex-floats}] Float placement, captions, subfloats, source notation.
	\item[\module{omnilatex-document}] Document metadata, custom \texttt{\textbackslash document*} commands.
	\item[\module{omnilatex-boxes}] Color boxes via tcolorbox.
	\item[\module{omnilatex-koma}] KOMA-Script sectioning, titlepage, TOC customization.
\end{description}

\section{Typography (\module{lib/typography/})}

\begin{description}
	\item[\module{omnilatex-fonts}] Font loading with graceful fallbacks.
	\item[\module{omnilatex-typesetting}] Line spacing, dashes, quotation marks.
	\item[\module{omnilatex-math}] Math typesetting, delimiters, physics macros, big-O notation.
	\item[\module{omnilatex-lists}] Custom list environments via enumitem.
\end{description}

\section{References (\module{lib/references/})}

\begin{description}
	\item[\module{omnilatex-hyperref}] Hyperlinks, PDF metadata, git SHA embedding.
	\item[\module{omnilatex-biblio}] Bibliography via biblatex (biber backend).
	\item[\module{omnilatex-glossary}] Glossaries, abbreviations, symbols via bib2gls.
\end{description}

\section{Language (\module{lib/language/})}

\begin{description}
	\item[\module{omnilatex-i18n}] Polyglossia integration, translations.
\end{description}

\section{Graphics (\module{lib/graphics/})}

\begin{description}
	\item[\module{omnilatex-graphics}] Image handling, SVG support (Inkscape).
	\item[\module{omnilatex-tikz-core}] TikZ/PGFPlots core, forest trees.
	\item[\module{omnilatex-tikz-engineering}] Engineering diagrams, P\&ID, flowcharts.
\end{description}

\section{Code (\module{lib/code/})}

\begin{description}
	\item[\module{omnilatex-listings}] Source code highlighting via minted.
\end{description}

\section{Tables (\module{lib/tables/})}

\begin{description}
	\item[\module{omnilatex-tables}] Table formatting (booktabs, tabularray, longtable).
\end{description}

\section{Utilities (\module{lib/utils/})}

\begin{description}
	\item[\module{omnilatex-utils}] TODO notes, git metadata, keyboard macros, censoring.
	\item[\module{omnilatex-colors}] Color definitions and schemes.
\end{description}

%% ─────────────────────────────────────────────────────
\chapter{Font System}
%% ─────────────────────────────────────────────────────

OmniLaTeX uses a three-font stack:

\begin{description}
	\item[Libertinus Serif + Math] Main text and mathematics. Always available via TeX Live.
	\item[Monaspace Neon] Monospace and code listings. Optional — falls back to Latin Modern Mono.
	\item[Atkinson Hyperlegible Next] Sans-serif elements. Optional — falls back to Libertinus Sans.
\end{description}

Missing optional fonts trigger a \texttt{\textbackslash ClassWarning} and degrade
gracefully. No manual configuration is needed.

%% ─────────────────────────────────────────────────────
\chapter{Build Modes}
%% ─────────────────────────────────────────────────────

OmniLaTeX supports three build modes, controlled via the \texttt{OMNILATEX\_MODE}
environment variable:

\begin{description}
	\item[dev] Development mode. Code listings use basic formatting, debug output enabled.
	\item[prod] Production mode (default). Full formatting, no debug output.
	\item[ultra] Ultra mode. Maximum quality, SVG conversion, no compromises.
\end{description}

%% ─────────────────────────────────────────────────────
\chapter{Reproducible Builds}
%% ─────────────────────────────────────────────────────

OmniLaTeX supports deterministic PDF generation via \texttt{SOURCE\_DATE\_EPOCH}:

\begin{minted}{bash}
SOURCE_DATE_EPOCH=$(git log -1 --format=%ct) latexmk -lualatex main.tex
\end{minted}

This sets all PDF timestamps to the git commit date, enabling byte-for-byte
reproducible builds.

%% ─────────────────────────────────────────────────────
\chapter{License}
%% ─────────────────────────────────────────────────────

OmniLaTeX is licensed under the Apache License 2.0.

\end{document}
